\chapter{Contexto y Estado del Arte}
\label{cap:estado}

\section{El problema: subtitular una clase para quien no puede seguirla}

\subsection{El marco normativo}

El derecho a la educación inclusiva de las personas con discapacidad está reconocido en el
artículo 24 de la Convención de Naciones Unidas sobre los Derechos de las Personas con
Discapacidad \cite{onu2006cdpd}, ratificada por España, y desarrollado en la legislación
española por el texto refundido de la Ley General de derechos de las personas con
discapacidad \cite{espana2013tr}, que obliga a los centros a realizar los ajustes
razonables necesarios. En la práctica, el ajuste habitual para el alumnado con discapacidad
auditiva en una clase magistral es humano: un intérprete de lengua de signos o un
transcriptor en directo.

Existe además una norma que fija qué es un subtítulo aceptable. La UNE
153010 \cite{aenor2012une153010} establece criterios de velocidad de lectura, segmentación y
presentación para el subtitulado dirigido a personas sordas. Este trabajo no la aplica como
requisito de conformidad, porque el sistema desarrollado no es un producto certificable; sí
la toma como referencia de dos decisiones de interfaz: que el texto se presente en unidades
sintácticas completas, y que el tamaño lo controle quien lee.

La distancia entre esa norma y lo que un reconocedor automático produce es precisamente el
objeto del trabajo. Y hay evidencia previa de que la métrica habitual no mide esa distancia:
al evaluar con personas sordas o con discapacidad auditiva la utilidad de subtítulos
generados automáticamente, la tasa de error de palabra predice mal la comprensión percibida,
porque no todos los errores estorban igual \cite{kafle2017captions}. Este trabajo llega a la
misma conclusión por una vía distinta, midiendo qué unidades léxicas sobreviven.

\subsection{Para quién funciona un subtítulo}
\label{sec:para-quien}

El colectivo que motiva la norma no es el único que se beneficia, y conviene documentarlo
porque amplía sustancialmente el alcance de un sistema como el que aquí se construye. La
revisión de más de cien estudios empíricos recogida por Gernsbacher
\cite{gernsbacher2015captions} concluye que el subtitulado en la misma lengua mejora la
comprensión, la atención y la memoria del material audiovisual en general, y que resulta
\emph{especialmente} útil en tres poblaciones: quienes son sordos o tienen discapacidad
auditiva, quienes están aprendiendo a leer, y quienes ven el material en una lengua que no
es la suya.

Sobre esa tercera población la evidencia es cuantitativa y antigua. La observación
fundacional es de Vanderplank \cite{vanderplank1988teletext}, que documentó cómo aprendices
de una lengua extranjera explotaban los subtítulos del teletexto para fijar vocabulario que
no habrían retenido solo escuchando. El metaanálisis de Montero Perez y colaboradores
\cite{monteroperez2013meta}, sobre dieciocho estudios, encuentra un efecto grande tanto en
comprensión oral como en adquisición de vocabulario, con un tamaño del efecto de 0,87 para
esta última. El mecanismo que se propone es el que cabría esperar: el subtítulo ancla la
cadena sonora a una forma escrita que el aprendiz sí reconoce, y convierte en léxico
conocido lo que de otro modo pasaría como ruido.

Ahora bien, esa literatura mide casi siempre \emph{vídeo grabado}, y este trabajo construye
un sistema \emph{en vivo}, que no es lo mismo: no hay posibilidad de pausar, ni de rebobinar,
ni de ver dos veces, que son precisamente las estrategias que el aprendiz emplea en los
estudios anteriores. El único trabajo localizado sobre transcripción en directo en un aula
de segunda lengua \cite{wang2023transcripts} obtiene un resultado más matizado, y merece la
pena reproducirlo con exactitud: las transcripciones en vivo \textbf{no} mejoraron las
calificaciones de ningún grupo, pero el alumnado de menor competencia lingüística completó un
24,6\% más de actividades de clase cuando disponía de ellas. Es decir, el efecto medido en
directo no aparece en el rendimiento sino en la \emph{participación}: el subtítulo no hace
que el alumno aprenda más deprisa, hace que pueda seguir la clase en lugar de descolgarse de
ella.

Esa distinción importa para no sobrevender el sistema. Lo que la evidencia sostiene es que
un subtítulo en vivo mantiene enganchado a quien de otro modo perdería el hilo; lo que no
sostiene todavía, porque no está medido, es que acelere el aprendizaje de la lengua en ese
escenario. Encaja además con lo que se observa en el conjunto del alumnado: en una encuesta
a 2.124 estudiantes de quince universidades \cite{linder2016captions}, el 98,6\% consideró
útiles los subtítulos y tres de cada cuatro los usaban como apoyo al estudio, sin tener
ninguna discapacidad ni cursar en lengua extranjera.

La conclusión que este trabajo extrae de esa literatura es de diseño y no de resultados: un
sistema de subtitulado de aula tiene una población objetivo bastante mayor que el colectivo
que lo justifica normativamente, y por tanto tiene sentido que la vista del alumnado sea
abierta y sin fricción, en lugar de un recurso que haya que solicitar. Es la decisión que se
describe en el capítulo~\ref{cap:desarrollo}, y esta sección es su fundamento.

Hay una consecuencia menos cómoda, y también se declara aquí. Si el sistema sirve a quien
está aprendiendo la lengua, entonces los modos de fallo estudiados en este trabajo le afectan
\emph{más} que al lector nativo, no menos: quien domina el castellano puede reconstruir por
contexto una negación perdida o un término mal transcrito, y quien lo está aprendiendo,
no. El argumento a favor de medir lo que el WER agregado esconde se refuerza, por tanto,
exactamente en la población para la que la evidencia de beneficio es más sólida.

\section{Reconocimiento automático del habla de propósito general}

El punto de partida de este trabajo es Whisper \cite{radford2023whisper}, un modelo
codificador-decodificador entrenado sobre una cantidad masiva de audio con supervisión
débil. Su relevancia para el problema planteado no está en alcanzar el mejor resultado en
un banco de pruebas concreto, sino en dos propiedades prácticas: funciona en español sin
necesidad de entrenamiento adicional, y es lo bastante robusto frente a condiciones de
grabación variables como para plantear su uso fuera de un estudio.

El modelo procesa el audio en ventanas de treinta segundos. Esa decisión de diseño, que
resulta natural para transcribir grabaciones completas, se convierte en un obstáculo cuando
se pretende subtitular en directo: el sistema no puede esperar medio minuto para mostrar la
primera palabra. La adaptación de un modelo pensado para procesamiento por lotes a un
escenario de flujo continuo es un problema abierto \cite{machacek2023whisperstreaming}, y
constituye una parte sustancial del trabajo de ingeniería descrito en el
capítulo~\ref{cap:desarrollo}.

\subsection{Una familia alternativa: los transductores}
\label{sec:transductores}

Whisper pertenece a la familia codificador-decodificador autorregresiva: genera el texto
palabra a palabra condicionando cada predicción en las anteriores. Existe una segunda
familia, el transductor neuronal \cite{graves2012transducer}, que emite símbolos alineados
con las tramas de audio y no dispone de ese bucle sobre el texto ya generado. Es la
arquitectura dominante en los sistemas pensados para funcionar en directo, precisamente
porque no necesita ver la señal completa para empezar a emitir.

Los sistemas recientes de esta familia combinan un codificador convolucional-atencional
\cite{gulati2020conformer, rekesh2023fastconformer} con una variante del transductor que
predice conjuntamente los símbolos y su duración \cite{xu2023tdt}, lo que reduce
drásticamente el número de pasos de decodificación y, con ello, el coste por segundo de
audio.

Esta sección existe porque el capítulo~\ref{cap:desarrollo} contrasta el modelo elegido con
un sistema de esta segunda familia. La diferencia arquitectónica no es un detalle
taxonómico: determina qué se puede controlar del proceso de decodificación, y ese control
resultó ser el criterio que decidió la comparación.

\section{Técnicas de adaptación al dominio}

Cuando un modelo general se aplica a un dominio concreto, la terminología específica y el
registro propio de ese dominio suelen degradar el resultado. Las estrategias para
corregirlo se ordenan por coste creciente.

\textbf{Condicionamiento por contexto.} La más barata consiste en aportar al decodificador
un texto previo que oriente sus predicciones, sin modificar el modelo. En Whisper esto se
implementa mediante un prompt que actúa como transcripción anterior. Es atractiva porque
no requiere entrenamiento ni datos anotados: bastaría con que el docente proporcionara el
glosario de la asignatura.

\textbf{Corrección posterior con un modelo de lenguaje.} Una segunda vía deja intacto el
reconocedor y repara su salida mediante un modelo de lenguaje, aprovechando que este
dispone de conocimiento léxico y sintáctico del que el reconocedor carece. La intuición es
sólida; los resultados de este trabajo la contradicen, y el capítulo~\ref{cap:desarrollo}
explica por qué.

\textbf{Ajuste fino con adaptadores.} La tercera modifica los pesos del modelo. El ajuste
completo resulta inviable en equipamiento modesto, pero LoRA \cite{hu2021lora} permite
entrenar únicamente unas matrices de rango bajo insertadas en las capas de atención,
reduciendo los parámetros entrenables en varios órdenes de magnitud y produciendo
adaptadores de unos pocos megabytes. Es lo que hace que esta técnica sea abordable en el
contexto de un trabajo como este.

\section{Recursos de habla en español}

La disponibilidad de corpus condicionó el diseño experimental más que ninguna otra
decisión, hasta el punto de que conviene documentar la búsqueda y no solo su resultado.

\subsection{El corpus objetivo}

El recurso que mejor se ajusta al problema es poliMedia, el repositorio de clases grabadas
de la Universitat Politècnica de València, del que el proyecto europeo transLectures
transcribió manualmente más de ciento quince horas \cite{deltoro2014translectures}. Reúne
las tres condiciones que este trabajo necesita: dominio educativo real, lengua española y
transcripción verificada por personas.

% AL ENVIAR LA SOLICITUD: sustituir por «Se cursó una solicitud institucional al grupo
% responsable, que no obtuvo respuesta dentro del plazo de este trabajo», y anotar la fecha
% en docs/solicitud-polimedia.md. Hasta entonces, esa frase seria falsa y comprobable.
No se encuentra publicado en ningún repositorio de acceso abierto, y su sitio web rechazó
la conexión de forma repetida durante la fase de auditoría de corpus, de modo que no fue
posible confirmar por esa vía ni el procedimiento de acceso ni la licencia. El repositorio
público del grupo distribuye otros conjuntos de datos, todos en inglés, pero no este: el
único camino es la solicitud directa, que queda redactada y dirigida al grupo responsable
como continuación inmediata del trabajo. Esta circunstancia es la limitación de fondo de
toda la parte experimental.

\subsection{Los recursos efectivamente empleados}

A falta del corpus objetivo se recurrió a sustitutos, escogidos para cubrir ejes distintos
del problema en lugar de acumular horas.

VoxPopuli \cite{wang2021voxpopuli} aporta discurso parlamentario en español peninsular.
Presenta dos propiedades que resultaron decisivas: incorpora una marca que identifica las
transcripciones verificadas manualmente, y anota la variedad dialectal del hablante. Es,
por eso, el único de los cinco sobre el que se realiza el ajuste fino.
FLEURS \cite{conneau2022fleurs} proporciona voz leída de dominio general, útil como techo
de calidad; conviene señalar que su configuración en español corresponde a la variedad
latinoamericana, sin equivalente peninsular.

Los dos recursos de habla espontánea proceden del proyecto CIEMPIESS de la Universidad
Nacional Autónoma de México \cite{hernandezmena2014ciempiess}: un conjunto de charlas
expositivas y otro de divulgación científica televisiva. Son los que más se acercan al
registro de una clase magistral entre los disponibles sin trámites, y también los de peor
calidad de referencia, como se explica a continuación. Se añadió por último un conjunto de
locución de medios en español peninsular \cite{kolobov2021mediaspeech}, que cubre el eje
dialectal a registro controlado.

Se descartó Common Voice \cite{ardila2020commonvoice} pese a su volumen: son frases leídas
y aisladas, sin la continuidad discursiva que caracteriza a una explicación de aula, de modo
que habría aportado horas sin aportar el fenómeno que se quería medir.

\subsection{La calidad de la referencia como variable}

Un hallazgo de la fase de búsqueda merece mención aquí porque afecta a la interpretación de
toda la literatura basada en estos recursos: la calidad de las transcripciones de
referencia varía sustancialmente entre corpus abiertos, y esa variación se confunde con
diferencias de rendimiento del modelo.

En los recursos de habla espontánea empleados, las referencias omiten la acentuación de
forma sistemática y contienen erratas. El efecto se cuantifica reevaluando el mismo
resultado con un normalizador que ignora las tildes: la columna \emph{$\Delta$ tildes} de la
tabla~\ref{tab:baseline-corpus} recoge esa diferencia. Sobre los corpus mexicanos alcanza
3,48 y 5,11 puntos, es decir, entre un cuarto y un tercio de todo el error atribuido al
modelo; sobre los peninsulares, con referencia verificada, cae a 0,07 puntos. La brecha no
mide reconocimiento: mide la referencia.

Se observaron además casos en que el sistema escribía correctamente una palabra que la
transcripción de referencia recogía mal, y era penalizado por ello: \emph{mallorca} donde
la referencia decía \emph{mayorca}, \emph{velázquez} donde decía \emph{velásquez}.

La consecuencia práctica quedó incorporada al procedimiento: antes de adoptar un corpus se
audita manualmente una muestra de sus referencias. Que un recurso disponga de transcripción
no implica que sirva como referencia.

\section{Evaluación de sistemas de reconocimiento}

\subsection{La métrica habitual y sus límites}

La tasa de error de palabra, definida como la distancia de edición entre referencia e
hipótesis normalizada por la longitud de la referencia, es la métrica estándar del campo.
Su fortaleza es la comparabilidad; su debilidad, que asigna el mismo coste a todas las
palabras y a todas las operaciones de edición.

Esa limitación está documentada en la literatura, que señala que la métrica opera en el
plano léxico sin consideración del contexto semántico, y que trata de forma idéntica errores
que destruyen la comprensión y errores irrelevantes para ella. Se han propuesto alternativas
que combinan la distancia de edición con señales de distancia semántica entre referencia e
hipótesis \cite{kim2021semdist}, y en el ámbito concreto del subtitulado accesible se ha
mostrado experimentalmente, con usuarios sordos o con discapacidad auditiva, que dos
transcripciones con el mismo WER pueden diferir mucho en cuánto se entiende de
ellas \cite{kafle2017captions}.

Conviene precisar también cómo se calcula aquí el intervalo de confianza de esa métrica. El
remuestreo por bootstrap \cite{efron1993bootstrap} es el procedimiento habitual para ello en
reconocimiento del habla, con una advertencia recogida en la literatura del campo desde hace
dos décadas \cite{bisani2004bootstrap}: la unidad de remuestreo debe ser el fragmento
y no la palabra, porque los errores se agrupan dentro de una misma locución y tratarlos como
independientes estrecha el intervalo de forma artificial.

Este trabajo llega a la misma conclusión por vía empírica y desde un dominio concreto: los
errores que perjudican a un alumno con discapacidad auditiva no son los que más pesan en la
métrica. Las mediciones que lo sustentan, y los instrumentos complementarios que se
proponen, se presentan en el capítulo~\ref{cap:desarrollo}.

\subsection{Alucinaciones}

Un modo de fallo específico de los modelos generativos aplicados al reconocimiento consiste
en producir texto fluido y verosímil sin correspondencia con el audio. A diferencia de un
error ordinario, resulta indetectable para el usuario final: la salida es gramatical y
coherente con el dominio.

Para un sistema de accesibilidad esta distinción es central. Un subtítulo ausente se
percibe como una carencia; uno inventado se lee como información. El trabajo incorpora por
ello la detección de salidas anómalas como criterio de selección de modelo, y no solo la
tasa de error.

\section{Subtitulado en tiempo real}

La literatura sobre adaptación de Whisper a escenarios de flujo continuo
\cite{machacek2023whisperstreaming} y la práctica documentada del sector coinciden en
varios puntos que sirvieron de referencia para el diseño de la aplicación: emplear ventanas
de entre cinco y diez segundos, solaparlas para no perder las palabras que caen en la
frontera, y apoyarse en detección de actividad de voz para decidir cuándo dar por terminada
una intervención. Para esto último la práctica habitual recurre a detectores entrenados
\cite{silero2021vad}; este trabajo emplea uno de energía, y la
sección~\ref{sec:limitaciones} razona por qué esa elección no invalida el hallazgo que
sostiene.

Coinciden también en desaconsejar el condicionamiento sobre la transcripción previa entre
ventanas, por el riesgo de propagación de errores. Este trabajo verificó esa recomendación
de forma independiente y obtuvo el mismo resultado, lo que constituye una validación
externa útil de la parte experimental.

La divergencia principal respecto a la práctica habitual está en el umbral de silencio
empleado para cerrar una intervención, y se discute en el capítulo~\ref{cap:desarrollo}
junto con la medición que la respalda.

\section{Posicionamiento del trabajo}

Este trabajo no propone una técnica nueva de adaptación ni un modelo nuevo. Su aportación
se sitúa en tres planos.

En el plano empírico, compara tres técnicas de adaptación sobre español con un diseño y un
criterio estadístico homogéneos, y contrasta su rendimiento con el de las decisiones de
ingeniería del sistema; una comparación que la literatura suele tratar por separado.

En el plano metodológico, documenta las condiciones bajo las cuales esas mediciones son
válidas, incluyendo varios casos en que un procedimiento aparentemente correcto produce
cifras verosímiles y falsas.

En el plano aplicado, construye un sistema desplegable en un centro educativo y lo emplea
como banco de pruebas, del que surgieron preguntas que la parte experimental no se habría
planteado.
